\section{Kalman filter 是动态系统上的 Gaussian 条件化}
\label{sec:13-Machine-Learning/07-Graphical-and-Sequential-Models/03}

假设你负责一个冷库的温度控制。墙上装着温度传感器，每十秒报一个数；你还有一个热平衡方程，描述了压缩机启动后温度大致怎样变化。问题是：传感器读数带着噪声，墙体蓄热和开门装货造成的扰动也没有写进方程。仅信传感器，噪声会让控制器反复误启动；仅信方程，未建模的扰动会慢慢把预测带偏。

HMM 让隐状态取有限类别，但温度是连续向量，漂移的传感器偏置也是连续量。把状态空间模型写成线性 Gaussian 形式：

\[
X_t=AX_{t-1}+Bu_t+W_t,
\qquad
W_t\sim\mathcal N(0,Q),
\]

\[
Y_t=CX_t+V_t,
\qquad
V_t\sim\mathcal N(0,R),
\]

其中过程噪声 \(W_t\)、观测噪声 \(V_t\) 与初始状态彼此独立。线性与 Gaussian 联手的后果是：所有 filtering posterior 永远保持 Gaussian。追踪状态的信念只需要两个量——均值与协方差——而它们的递推公式有闭式。

冷库的传感器偏置尤其值得警惕。若真实读数满足

\[
Z_t=T_t+b_t+\varepsilon_t,
\]

而模型把偏置 \(b_t\) 硬编码为零，多收数据只会让 filter 对错误温度越来越确信。把偏置扩入状态 \(b_t=b_{t-1}+\zeta_t\) 是正确方向，但温度与偏置能否从现有观测中分离，取决于下文的可观测性，不取决于递推步数。

\subsection{预测：把上一时刻的信念推过动力学}

设已有

\[
X_{t-1}\mid y_{1:t-1}
\sim
\mathcal N(m_{t-1},P_{t-1}).
\]

线性变换取走确定性部分，过程噪声独立加入：

\[
m_t^-=Am_{t-1}+Bu_t,
\]

\[
P_t^-=AP_{t-1}A^{\trans}+Q.
\]

上标 \(-\) 表示还没看到 \(y_t\)。即使均值沿确定性动力学一步不差，\(Q\) 仍会撑开协方差。若代码里漏掉 \(Q\)，filter 将对自身估计越来越过度自信，最终把真实观测当成异常值拒绝。

这两个公式在结构上就是 HMM 前向算法在连续 Gaussian 状态的版本，并无新概念。真正的区分点在下文。

\subsection{更新：联合 Gaussian 的条件公式}

新观测到达时，预测状态与观测构成联合 Gaussian 分布：

\[
\begin{bmatrix}
X_t\\Y_t
\end{bmatrix}
\sim
\mathcal N\left(
\begin{bmatrix}
m_t^-\\Cm_t^-
\end{bmatrix},
\begin{bmatrix}
P_t^- & P_t^-C^{\trans}\\
CP_t^- & CP_t^-C^{\trans}+R
\end{bmatrix}
\right).
\]

记 innovation（观测与预测之差）

\[
e_t=y_t-Cm_t^-,
\]

innovation covariance

\[
S_t=CP_t^-C^{\trans}+R.
\]

多元 Gaussian 的条件公式给出 posterior：

\[
K_t=P_t^-C^{\trans}S_t^{-1},
\]

\[
m_t=m_t^-+K_te_t,
\]

\[
P_t=P_t^--K_tCP_t^-.
\]

\(K_t\) 的分母是观测噪声的协方差，分子是预测状态与观测的协方差。因此，观测噪声 \(R\) 大时更新量小——传感器不可靠就不要太信它。预测不确定性 \(P_t^-\) 大且观测能看到该方向时，更新量大——我自己都不确定，你告诉我什么我就多听。

这就是 Kalman filter 的核心计算元件。没有新概率工具，就是把 Gaussian 的条件公式按时间顺序反复调用了一次。

\subsection{不可观测方向不会被数据凭空修复}

某个状态方向长期不影响 \(CX_t\) 时，观测无法直接约束它。动力学 \(A\) 可能逐步把该方向耦合到可测方向，也可能让它永远藏在不可观测子空间里。线性系统的 observability 矩阵

\[
\mathcal O=
\begin{bmatrix}
C\\CA\\ \vdots\\ CA^{d-1}
\end{bmatrix}
\]

回答了这个问题：若 \(\mathcal O\) 满列秩，有限时间内的无噪观测足以区分状态；若不满秩，再小的 \(R\) 也不能从数据恢复不可观测分量。

这解释了为什么 posterior covariance 不只是噪声大小的函数，还反映模型结构允许观测约束哪些方向。回到冷库例子：如果传感器只测温度而不反映偏置，且动力学不把偏置耦合到温度上，偏置就是对所有数据视而不见的幽灵变量——即使连续读数一整天，filter 的偏置估计方差也不会缩小。

\subsection{filtering、prediction 与 smoothing 回答不同问题}

三种查询的可用信息窗口不同：

Filtering 计算 \(p(x_t\mid y_{1:t})\)，只用当前及过去观测。Prediction 计算 \(p(x_{t+h}\mid y_{1:t})\)，推向未来。Smoothing 在整段数据到达后计算 \(p(x_t\mid y_{1:T})\)，\(T>t\)，允许未来观测回头修正过去状态。

离线轨迹重建可以用 smoothing；实时控制不能在时刻 \(t\) 使用尚未发生的 \(y_{t+1:T}\)。评价时若误将平滑状态当作在线特征，就发生了时间泄漏——与在训练时偷看测试集是同一种结构错误。

\subsection{数值实现应解线性系统而不是求逆}

公式写 \(S_t^{-1}\) 是为了表达线性算子的语义。计算时用 Cholesky 分解解 \(S_tZ=CP_t^-\) 等效的线性系统，而不是显式求逆矩阵。协方差更新的简式 \(P_t=P_t^--K_tCP_t^-\) 在舍入误差积累下可能失去对称性或半正定。更稳定的 Joseph form

\[
P_t
=(I-K_tC)P_t^-(I-K_tC)^{\trans}
+K_tRK_t^{\trans}
\]

以平方形式组合两项正半定矩阵，数值上更抗扰动。平方根 filter 进一步直接传播 Cholesky 因子，在病态问题中可以多撑几个数量级。这些不是附属的工程细节：负方差或非对称协方差会破坏下一次 Gaussian 条件化的前提，让整个递推变得无意义。

\subsection{非线性与非 Gaussian 击穿闭式递推}

若动力学或观测是非线性的，Gaussian 经过非线性变换后不再 Gaussian。Extended Kalman filter 在当前均值处线性化，误差依赖局部 Jacobian 和不确定性范围；unscented 方法用确定性 sigma points 近似矩传播；particle filter 则放弃 Gaussian 近似，用带权样本表示任意形状 posterior。

这些方法不是在给同一套闭式公式做数值加速。它们在闭式不再成立后，各自选择了不同近似原则。多峰 posterior 用单个均值和协方差描述，会把概率质量平均到实际上没有样本的低概率区域——此时即使线性化再精确，也无法恢复被 Gaussian 族丢掉的模态结构。

Kalman filter 的力量来自一组非常具体的闭包条件：线性映射、Gaussian 噪声与 Markov 状态转移。看清这些条件，也就看清了什么时候应继续信任它，什么时候必须换模型结构而不是继续微调协方差矩阵的对角元。
